\section{Local Variable Names}\label{sec:NamesForLocalVariables}
The identifiers for local variables have to be in mixed case, with a lowercase first letter. Internal words start with capital letters\footnote{see chapter \tqvref{sec:CaseOfNames}}. The names should contain neither the underscore (“\verb#_#”) nor the dollar sign (“\verb#$#”) characters, even though both are allowed by the language specification; they also should not start with those characters. Digits are forbidden as first characters already by the language specification, but perfectly allowed as additional characters. But they should be used with care and only were necessary.

Variable names should be short yet meaningful – with a clear priority on being meaningful. The choice of a variable name should be mnemonic – that is, designed to indicate the intent of its use to the casual observer. As for class names, abbreviations and acronyms should be avoided. Using the digit '\verb#4#' as replacement for the word “\verb#for#”, and '\verb#2#' for the word “\verb#to#” is acceptable, yet not really recommended.

In case the first part of the variable name has to be an acronym, it will be written with all lower case:
\begin{lstlisting}
String htmlHeader; // OK
\end{lstlisting}
instead of
\begin{lstlisting}
String hTMLHeader; // AVOID!
\end{lstlisting}

Single-letter variable names should be avoided except for temporary “throwaway” variables. Some of these are widely used and have already a fixed meaning; so is \lstinline|i| very common for the run value of a classical \lstinline|for| loop:
\begin{lstlisting}
for( var i = 0; i < max; ++i )
{
    …
}
\end{lstlisting}

\lstinline|e| is the common name for the exception in a \lstinline|catch| block:
\begin{lstlisting}
try
{
    …
}
catch( final IOException e )
{
    //---* Handle the exception *------------------------------------
    …
}    
\end{lstlisting}

For temporary strings, \lstinline|s| is common, as \lstinline|o| is for (temporary) objects of an unspecified type.

\label{lbl:RetValue}Some names are fixed: the name of the value that is returned from a method has to be \lstinline|retValue| if the method does not return a field or attribute or \lstinline|this|:
\begin{lstlisting}
// OK
public final int getValue() { return m_Value; }

// AVOID!! Use 'retValue' here instead of 'i'!!
public final int signum( double d )
{
    final var i = (d < 0.0) ? -1 : ((d > 0.0) ? 1 : 0);
    
    //---* Done *----------------------------------------------------
    return i;
}   //  signum()

// ACCEPTABLE (BARELY)
public final int signum( double d )
{
    return (d < 0.0) ? -1 : ((d > 0.0) ? 1 : 0);
}   //  signum()

// RECOMMENDED
public final int signum( double d )
{
    final var retValue = 0;
    if( d < 0.0 )
    {
        retValue = -1
    }
    else
    {
        retValue = 1;
    }

    //---* Done *----------------------------------------------------
    return retValue;
}   //  signum()

public final int signum( double d )
{
    final var retValue = (d < 0.0) ? -1 : ((d > 0.0) ? 1 : 0);

    //---* Done *----------------------------------------------------
    return retValue;
}   //  signum()
\end{lstlisting}

In the same way, \lstinline|result| is reserved as the return value for lambdas (see chapter \tqfullvref{sec:LambdaResults}) and for the results of \lstinline|case| blocks (see chapter \tqfullvref{sec:CaseResults}). The name \lstinline|result| should not be used for a local variable in any other context.

A temporary buffer is named \lstinline|buffer|, no matter if it is a \lstinline|java.lang.StringBuilder| or \lstinline|java.lang.StringBuffer|, an implementation of \lstinline|java.util.List|, another collection type, or any other type that can be used as a buffer.

Do not use the names \lstinline|in|, \lstinline|out|, and \lstinline|err| for your local variables. Although these names are not reserved words, they should be treated as such, as that allows to import the default streams from the class \lstinline|java.lang.System|\autocite{ORACLE_DOC_SYSTEM_CLASS} statically.

Usually you would use the singular form for a variable name, but for collections or arrays, the plural form can be perfectly correct:
\begin{lstlisting}
final Collection<Component> values = m_Components.values();
for( final Component value : values )
{
    …
}    
\end{lstlisting}

Using the class name (with a lower case first letter, of course) for a variable name is completely acceptable in case this is sufficient to explain the function of the variable; but the prefix “\verb#my#” or “\verb#the#” has to be avoided.
\begin{lstlisting}
// RECOMMENDED
public final String retrieveMessage( final String messageKey, 
    final String... additionalInfo )
{
    final MessageProvider messageProvider = getMessageProvider();
    final String rawMessage = 
        messageProvider.getMessage( messageKey );
    final var retValue = String.format( rawMessage, additionalInfo );
    
    //---* Done *----------------------------------------------------
    return retValue;
}   //  retrieveMessage()

// AVOID!!!! Don't use "my" or "the" prefix!!!
public final String retrieveMessage( final String string, 
    final String... strings )
{
    final MessageProvider theMessageProvider = getMessageProvider();
    final String myString = theMessageProvider.getMessage( string );
    final var retValue = String.format( myString, strings );
    
    //---* Done *----------------------------------------------------
    return retValue;
}   //  retrieveMessage()
\end{lstlisting}

Sometimes it is also a good idea to use the class name of the type as a suffix or prefix for the name, especially if the current block declares more than one variable of that type. In such a case, the rest of the name could be used to indicate how variables belong together.
\begin{lstlisting}
// RECOMMENDED
public final ResultData callSpecialService()
{
    SpecialClientAgent clientAgent = obtainClientAgent();
    …
}   //  callSpecialService()

// AVOID!!!! Don't use "my..." prefix!!! And don't abbreviate!
public final ResultData callSpecialService()
{
    SpecialClientAgent myCA = obtainClientAgent();
    …
}   //  callSpecialService()

// RECOMMENDED
public final ResultData loadData( Connection connection, … )
{
    …
    final Statement customerStatement = …
    final Statement orderStatement = …

    …

    final ResultSet customerResultSet = customerStatement.execute();
    …

    final ResultSet orderResultSet = orderStatement.execute();
    …
}   //  loadData()
\end{lstlisting}

Do not (never!) use a prefix to indicate the type of a variable (Hungarian Notation\autocite{WIKIPEDIA:Hungarian_notation,Jack:HungarianNotation}\index{Hungarian Notation})! Java is a strongly typed language and the compiler will take care that variables are only used according to their types.\footnote{The use of the Hungarian Notation\index{Hungarian Notation} – in particular the ‘Systems Hungarian’ variant – for the naming of variables is discouraged for strongly typed languages like Java or C++\index{C++}, because it does not provide any benefit.\newline For languages like PHP\index{PHP}, C\index{C}, JavaScript\index{JavaScript} or also Groovy\index{Groovy}, this may be different: for these languages the type prefix may help the programmers to assign the semantically correct type to a variable when syntactically any (or most) types are correct.} The use of the \lstinline|var| keyword\footnote{\lstinline|var| is in fact not a keyword, but a special type.} in Java programs does not weaken the strong typing of Java in any way – refer to chapter \tqfullvref{sec:LocalVariableTypeInference}.

Local variables inside lambdas are just like any other local variable and their naming follows the same rules. The only exception is that a lambda does not use \lstinline|retValue| for a return value, but \lstinline|result| – refer to chapter \tqfullvref{sec:LambdaResults} for the details.

%==============================================================================
% End of file!!
%==============================================================================
